% ============================================================
%  ON THE CATEGORICAL MECHANICS OF REAL-TIME EVENT SELECTION
%  IN BOUNDED PHASE SPACE: COMPOSITION-INFLATION, TEMPORAL
%  PROGRAMMING, AND S-ENTROPY TRAJECTORY COMPLETION
%  IN LHC TRIGGER SYSTEMS
% ============================================================

\documentclass[twocolumn,10pt,a4paper]{article}

\usepackage[utf8]{inputenc}
\usepackage[T1]{fontenc}
\usepackage{lmodern}
\usepackage[a4paper,top=2cm,bottom=2.2cm,left=1.8cm,right=1.8cm,columnsep=0.6cm]{geometry}
\usepackage{amsmath,amssymb,amsthm,mathtools}
\usepackage{bm}
\usepackage{booktabs}
\usepackage{array}
\usepackage{enumitem}
\usepackage{graphicx}
\usepackage{float}
\usepackage{xcolor}
\usepackage{microtype}
\usepackage[round,sort&compress]{natbib}
\usepackage[colorlinks=true,linkcolor=blue!60!black,
            citecolor=blue!60!black,urlcolor=blue!60!black]{hyperref}
\usepackage{fancyhdr}

\pagestyle{fancy}
\fancyhf{}
\fancyhead[LE,RO]{\thepage}
\fancyhead[RE]{\small Categorical Mechanics of LHC Trigger Systems}
\fancyhead[LO]{\small Sachikonye}
\renewcommand{\headrulewidth}{0.4pt}

% ---- Theorem environments ------------------------------------------------
\newtheoremstyle{thmstyle}{\topsep}{\topsep}{\itshape}{}{\bfseries}{.}{.5em}{}
\theoremstyle{thmstyle}
\newtheorem{theorem}{Theorem}[section]
\newtheorem{lemma}[theorem]{Lemma}
\newtheorem{proposition}[theorem]{Proposition}
\newtheorem{corollary}[theorem]{Corollary}
\theoremstyle{definition}
\newtheorem{definition}[theorem]{Definition}
\newtheorem{axiom}{Axiom}
\theoremstyle{remark}
\newtheorem{remark}[theorem]{Remark}

% ---- Custom commands -----------------------------------------------------
\newcommand{\tP}{t_{\mathrm{P}}}
\newcommand{\lP}{\ell_{\mathrm{P}}}
\newcommand{\nP}{n_{\mathrm{P}}}
\newcommand{\kB}{k_{\mathrm{B}}}
\newcommand{\R}{\mathbb{R}}
\newcommand{\N}{\mathbb{N}}
\newcommand{\Zp}{\mathbb{Z}^{+}}
\newcommand{\Parts}{\mathcal{P}}
\newcommand{\Sfl}{S_{\flat}}
\newcommand{\Sk}{S_{k}}
\newcommand{\St}{S_{t}}
\newcommand{\Se}{S_{e}}
\newcommand{\dP}{\Delta P}
\newcommand{\Trig}{\mathcal{T}}
\DeclareMathOperator{\catpow}{\kappa}

% =========================================================
\title{\textbf{On the Categorical Mechanics of Real-Time\\
Event Selection in Bounded Phase Space:\\
Composition-Inflation, Temporal Programming,\\
and S-Entropy Trajectory Completion\\
in LHC Trigger Systems}}

\author{
Kundai Farai Sachikonye\\
Technical University of Munich\\
Chair of Proteomics and Bioanalytics\\
\texttt{kundai.sachikonye@wzw.tum.de}
}
\date{\today}
% =========================================================

\begin{document}
\maketitle

% =========================================================
\begin{abstract}
\noindent
We develop a unified mathematical framework for real-time event selection
in Large Hadron Collider (LHC) trigger systems, grounded in three
independently motivated but mutually reinforcing structures: (i) the
\emph{composition-inflation mechanism}, which establishes that a trigger
system with $d$ independent detector subsystems observing $n$ bunch
crossings can discriminate $T(n,d) = d(d+1)^{n-1}$ categorically
distinguishable event patterns---an exponential, not linear, state
count; (ii) \emph{temporal programming}, a paradigm in which the sole
runtime datum is the timing deviation $\dP(k) = T_{\mathrm{ref}}(k) -
t_{\mathrm{rec}}(k)$ of the $k$-th detector signal from the LHC bunch
clock, and in which the trigger's cell registry constitutes a formally
verified, structurally incorruptible specification with no injection
surface; and (iii) \emph{S-entropy backward trajectory completion}, in
which the trigger's classification task is identified with navigation
from a known endpoint (the detected event) to its root (the originating
physics process) in a ternary refinement hierarchy, achieving
$\mathcal{O}(\log_3 N)$ complexity via virtual sub-state decompositions.

We derive the \emph{Planck depth} $\nP = 60$ for the LHC bunch clock
($\nu_B = 40.079\,\mathrm{MHz}$, $d = 3$): after 60 bunch crossings
($1.496\,\mu\mathrm{s}$), the trigger's categorical pattern count
$T(60,3) = 3 \cdot 4^{59} \approx 7.3 \times 10^{35}$ exceeds the
ratio $\tau_B/\tP = 4.63 \times 10^{35}$ of the bunch-crossing period
to the Planck time. The current L1 trigger latency of 100 bunch
crossings operates at depth $n/\nP \approx 1.7$, placing it entirely
within the super-Planck categorical resolution regime.

We establish a formal identification between \emph{virtual particles}
in quantum field theory and \emph{virtual sub-states} in the S-entropy
framework: off-shell propagation ($q^2 \neq m^2$) corresponds
precisely to sub-coordinate decompositions outside the physical unit
cube ($\mathbf{S} \notin [0,1]^3$), and the five canonical properties
of virtual particles---enabling transitions, zero net work, path
opacity, mediating catalytic power, and necessity for exponential
speedup---are theorems of the categorical framework, not physical
assumptions. The Ward--Takahashi identities of gauge field theory are
identified with the Local-Global Decoupling Theorem, and S-matrix
observability is identified with Path Opacity.

We derive all seven fundamental constants in categorical angular units:
$c = 2\pi\,\mathrm{rad/tick}$, $\hbar = E_{\mathrm{tick}}/(2\pi)$,
$\kB = E_{\mathrm{tick}}/\ln(d+1)$, $m_0 = E_{\mathrm{tick}} \nu_0
/(4\pi^2)$, with the gravitational constant $G$ remaining the unique
irreducible coupling not reducible to tick counting and $2\pi$ geometry.
The HL-LHC upgrade to $d = 4$ independent subsystems reduces the
Planck depth from 60 to 52, recovering eight bunch crossings of
categorical resolution headroom.

\bigskip
\noindent\textbf{Keywords:} LHC trigger, composition-inflation,
temporal programming, S-entropy, backward trajectory completion,
virtual sub-states, virtual particles, Planck depth, Ward identity,
S-matrix, categorical mechanics, bunch crossing, coincidence detection
\end{abstract}

\tableofcontents

% =========================================================
\section{Introduction}
\label{sec:introduction}
% =========================================================

\subsection{The Event Selection Problem}

The Large Hadron Collider at CERN produces proton--proton collisions at
a rate of 40.079\,MHz~\citep{evans2008lhc,lhc_bunch_structure}. Each
collision, if fully recorded, generates approximately 1--2\,MB of raw
detector data~\citep{atlas2008detector,cms2008detector}. The aggregate
data rate---approaching 80\,TB per second---cannot be recorded by any
existing storage technology. The trigger system must reduce this rate
to approximately 1\,kHz for permanent storage, rejecting
99.997\,\% of all events in real time~\citep{atlas_trigger_tdr,
atlas_trigger_run2}.

This event selection problem has three structural properties that
distinguish it from ordinary digital filtering:

\begin{enumerate}[leftmargin=1.5em,itemsep=0.2em]
\item \textbf{Real-time irreversibility.} The decision to discard an
event is permanent: the detector data is overwritten in the analog
pipeline before the next bunch crossing arrives. No second chance is
available~\citep{atlas_hlt2003,cms_trigger2002}.

\item \textbf{Timing primacy.} The sole physically guaranteed
distinction between a signal event and background is the timing of
detector signals relative to the LHC bunch clock. Content-bearing
information (particle identities, masses, charges) is inferred from
timing patterns, not transmitted as content~\citep{leo1994techniques}.

\item \textbf{Exponential state space.} The number of distinct event
topologies accessible at each bunch crossing is not $n$ (the number
of crossings observed) but grows exponentially with $n$---a fact
not captured by any existing formal framework for trigger design.
\end{enumerate}

The present paper establishes a rigorous mathematical framework that
makes these three properties precise, derives their consequences, and
produces new design principles for the LHC and its HL-LHC upgrade.

\subsection{Overview of Main Results}

The central results are organised as follows.

\textbf{Composition-Inflation (Section~\ref{sec:inflation}).}
We prove that a trigger system with $d$ independent detector
subsystems, observing $n$ bunch crossings, can discriminate
\begin{equation}
T(n,d) = d \cdot (d+1)^{n-1}
\end{equation}
categorically distinguishable event patterns. This exponential growth
arises from two independent multiplications: $2^{n-1}$ ordered
compositions of $n$ bunch crossings (the temporal structure), and $d^k$
dimensional labelings of $k$-part compositions (the spatial structure).
The na\"ive linear count of $n$ states per channel misses the
combinatorial richness by a factor $d(d+1)^{n-1}/n$, which exceeds
$10^{28}$ for $n = 56$ at $d = 3$.

\textbf{Planck Depth for the LHC (Section~\ref{sec:planck_depth}).}
We define the \emph{Planck depth} $\nP$ as the minimum number of bunch
crossings at which $T(\nP, d)$ exceeds the ratio $\tau_B / \tP$ of
the bunch-crossing period to the Planck time. For the LHC at $d = 3$,
$\nP = 60$ (1.496\,$\mu$s). The current L1 trigger latency of 100
bunch crossings operates at depth 1.7 times the Planck depth.

\textbf{Temporal Programming (Section~\ref{sec:temporal}).}
We formalise the trigger as a \emph{Tempus program}: a compiled cell
registry over the timing-deviation space $\dP(k) = T_{\mathrm{ref}}(k)
- t_{\mathrm{rec}}(k)$. The program has no content parser and no
injection surface, satisfying a \emph{structural incorruptibility}
theorem: no signal whose timing does not originate from the LHC bunch
crossing can satisfy any trigger cell condition.

\textbf{Backward Trajectory Completion (Section~\ref{sec:backward}).}
We identify the trigger decision problem as backward navigation in a
ternary refinement hierarchy: the endpoint (detector hit pattern) is
known, and the root (originating physics process) must be found. Using
the S-entropy embedding of the hierarchy into $[0,1]^3$ with the
product Fisher metric, backward navigation achieves
$\mathcal{O}(\log_3 N)$ complexity. Restricting to physical
intermediate states collapses to $\Theta(N)$, proving that
\emph{virtual sub-states are computationally necessary}.

\textbf{Virtual Particles as Virtual Sub-States (Section~\ref{sec:virtual}).}
We prove a formal identification theorem: virtual particles in quantum
field theory are virtual sub-states in the S-entropy framework. The
off-shell condition ($q^2 \neq m^2 c^4$) corresponds to lying outside
$[0,1]^3$; the Feynman propagator corresponds to the mean-recovery
constraint; gauge invariance corresponds to Local-Global Decoupling;
and S-matrix observability corresponds to Path Opacity.

\textbf{Constant Reduction (Section~\ref{sec:constants}).}
All seven fundamental constants reduce to tick counts, $2\pi$ geometry,
and tick energy in the categorical angular framework, with $G$ as the
unique irreducible coupling. This implies the Planck time is not a
fundamental resolution limit but a composition depth: $\tP
\leftrightarrow \nP \approx 56$--$72$ for any physical oscillator.

% =========================================================
\section{Mathematical Foundations: Bounded Phase Space}
\label{sec:foundations}
% =========================================================

\subsection{The Bounded Phase Space Axiom}

All physical systems in this framework occupy bounded regions of phase
space. This is not an assumption but a physical necessity: no system
possesses infinite energy or infinite momentum.

\begin{axiom}[Bounded Phase Space]
\label{ax:bounded}
Every physical system occupies a bounded region $\Omega \subset
\mathcal{P}$ of phase space with finite Liouville measure:
$\mu(\Omega) < \infty$.
\end{axiom}

\begin{theorem}[Partition Existence \citealp{arnold1989methods}]
\label{thm:partition-existence}
For any bounded region $\Omega$ with $\mu(\Omega) < \infty$, there
exists a finite partition $\{\Omega_i\}_{i=1}^N$ with $N < \infty$
such that $\Omega = \bigsqcup_{i=1}^N \Omega_i$.
\end{theorem}

\begin{proof}
By compactness of bounded sets in phase space (Heine--Borel theorem),
any open cover of $\Omega$ admits a finite subcover. The partition
follows from refinement of such a cover; finiteness of $N$ follows
from $\mu(\Omega) < \infty$ requiring each $\Omega_i$ to have positive
measure. \qed
\end{proof}

\subsection{Oscillatory Necessity and Action-Angle Variables}

\begin{theorem}[Poincar\'e Recurrence \citealp{poincare1890trois}]
\label{thm:poincare}
Let $\Sigma_E = H^{-1}(E)$ be a compact energy surface with
$\mu(\Sigma_E) < \infty$. For any measurable set $A \subset \Sigma_E$
with $\mu(A) > 0$, almost every trajectory starting in $A$ returns to
$A$ infinitely often.
\end{theorem}

\begin{corollary}[Oscillatory Necessity]
\label{cor:oscillatory}
Every compact Hamiltonian system without fixed points is oscillatory:
no coordinate can undergo monotonic escape, and every trajectory returns
to any neighbourhood of its initial condition.
\end{corollary}

\begin{proof}
Monotonic escape contradicts Theorem~\ref{thm:poincare} on a compact
energy surface. Absence of fixed points (guaranteed for systems at
$T > 0$ by the Third Law) forces the trajectory to oscillate. \qed
\end{proof}

For integrable systems, the Arnold--Liouville theorem~\citep{arnold1989methods}
guarantees the existence of action-angle variables $(J_i, \theta_i)$
with quasi-periodic motion $\dot\theta_i = \omega_i(J)$. Closure of
orbits requires $n_i \omega_i = n_j \omega_j$ for positive integers
$n_i$; the tuple $(n_1, n_2, n_3)$ in three dimensions labels the
\emph{partition state}.

\subsection{Partition Coordinates and Capacity}

\begin{definition}[Partition State]
A \emph{partition state} of a bounded three-dimensional Hamiltonian
system is an equivalence class of phase-space points sharing the same
closure label $(n, \ell, m, s)$ where:
$n \geq 1$ (principal level), $0 \leq \ell < n$ (angular sublevel),
$-\ell \leq m \leq \ell$ (azimuthal projection), $s = \pm\tfrac12$
(rotational parity). The set of all partition states is $\Parts$.
\end{definition}

\begin{theorem}[Partition Capacity]
\label{thm:capacity}
The number of distinct partition states at principal level $n$ is
$C(n) = 2n^2$.
\end{theorem}

\begin{proof}
At level $n$, $\ell$ ranges over $\{0,1,\ldots,n-1\}$; for each $\ell$,
$m$ ranges over $2\ell + 1$ values; and $s$ contributes a factor of 2:
\[
C(n) = 2\sum_{\ell=0}^{n-1}(2\ell+1) = 2\left(n + 2\cdot\tfrac{n(n-1)}{2}\right) = 2n^2.\qed
\]
\end{proof}

The cumulative count is $N_{\mathrm{state}}(n) = \tfrac{n(n+1)(2n+1)}{3}$,
growing cubically.

% =========================================================
\section{Composition-Inflation}
\label{sec:inflation}
% =========================================================

\subsection{Integer Compositions}

\begin{definition}[Composition]
A \emph{composition} of a positive integer $n$ is an ordered sequence
$(c_1, c_2, \ldots, c_k)$ of positive integers summing to $n$. Two
compositions differing in order are distinct.
\end{definition}

\begin{theorem}[Composition Count \citealp{andrews1976partitions}]
\label{thm:comp-count}
The total number of compositions of $n$ is $2^{n-1}$.
\end{theorem}

\begin{proof}
Place $n$ objects in a row; there are $n-1$ gaps between consecutive
objects. A composition corresponds to a binary choice (place/do not
place a divider) at each gap: $2^{n-1}$ configurations in bijection
with compositions. \qed
\end{proof}

The number of compositions into exactly $k$ parts is
$\binom{n-1}{k-1}$ (choose $k-1$ of the $n-1$ gaps for dividers).

\subsection{Dimensional Labeling}

The S-entropy space is the three-dimensional unit cube $[0,1]^3$ with
coordinates $(\Sk, \St, \Se)$ representing knowledge, temporal, and
categorical entropy dimensions, respectively.

\begin{definition}[Labeled Composition]
A \emph{labeled composition} of $n$ in $d$ dimensions is a pair
$(\mathbf{C}, \mathbf{L})$ where $\mathbf{C} = (c_1,\ldots,c_k)$ is
a composition and $\mathbf{L} = (\ell_1,\ldots,\ell_k)$ with
$\ell_i \in \{1,\ldots,d\}$ specifies the entropy dimension being
refined in the $i$-th part.
\end{definition}

\begin{theorem}[Composition-Inflation Formula]
\label{thm:inflation}
The total number of labeled compositions of $n$ in $d$ dimensions is
\begin{equation}
\label{eq:inflation}
T(n,d) = d \cdot (d+1)^{n-1}.
\end{equation}
\end{theorem}

\begin{proof}
Summing over all part counts:
\begin{align}
T(n,d) &= \sum_{k=1}^{n}\binom{n-1}{k-1}d^k
= d\sum_{j=0}^{n-1}\binom{n-1}{j}d^{j}
= d(1+d)^{n-1},
\end{align}
where the last step is the binomial theorem with $x = d$, $m = n-1$. \qed
\end{proof}

\begin{corollary}[S-Entropy Space]
For $d = 3$: $T(n,3) = 3 \cdot 4^{n-1}$. After $n = 56$ cycles, $T
= 3 \cdot 4^{55} \approx 3.8 \times 10^{33}$, exceeding $10^{33}$
distinguishable trajectories.
\end{corollary}

\begin{proposition}[Growth Properties]
\label{prop:growth}
The labeled composition count satisfies:
\begin{enumerate}[nosep,label=(\roman*)]
\item $T(1,d) = d$;
\item $T(n+1,d) = (d+1)\cdot T(n,d)$ (each tick multiplies states by $d+1$);
\item $\log T(n,d) = \log d + (n-1)\log(d+1)$;
\item $T(n,d)/n \to \infty$ as $n \to \infty$ (superlinear growth).
\end{enumerate}
\end{proposition}

\begin{proof}
Direct from Theorem~\ref{thm:inflation}. \qed
\end{proof}

\subsection{Angular Resolution and the Planck Constraint}

\begin{definition}[Angular Resolution]
For $T(n,d)$ distinguishable trajectories distributed over the phase
circle $[0,2\pi)$:
\[
\Delta\theta(n,d) = \frac{2\pi}{T(n,d)} = \frac{2\pi}{d(d+1)^{n-1}}.
\]
\end{definition}

\begin{theorem}[No Planck Limit on Angular Resolution]
\label{thm:no-planck}
The angular resolution $\Delta\theta(n,d)$ is dimensionless (units of
radians, not seconds) and therefore not subject to the Planck time
constraint. For any $\epsilon > 0$, there exists $n_0 = 1 +
\lceil\log_{d+1}(2\pi/(d\epsilon))\rceil$ such that $\Delta\theta <
\epsilon$ for all $n \geq n_0$.
\end{theorem}

\begin{proof}
The Planck time $\tP = \sqrt{\hbar G/c^5}$ has dimension
$[\tP] = \mathrm{s}$. The Heisenberg bound $\Delta E \cdot \Delta t
\geq \hbar/2$ constrains energy--time products; both factors carry
physical dimensions. The angular resolution $\Delta\theta =
2\pi/T(n,d)$ is the ratio of two dimensionless quantities and carries
no time dimension: $[\Delta\theta] = \mathrm{rad} = 1$. Neither the
Heisenberg argument nor the gravitational collapse argument applies to
a dimensionless ratio. The lower bound on $n_0$ follows from
$d(d+1)^{n-1} > 2\pi/\epsilon$, solved for $n$. \qed
\end{proof}

% =========================================================
\section{The Planck Depth for LHC Trigger Systems}
\label{sec:planck_depth}
% =========================================================

\begin{definition}[Planck Depth]
\label{def:planck-depth}
The \emph{Planck depth} $\nP$ for a reference oscillator with period
$\tau_{\mathrm{osc}}$ in $d$-dimensional S-entropy space is the minimum
cycle count for which
\[
T(\nP, d) \geq \frac{\tau_{\mathrm{osc}}}{\tP}.
\]
\end{definition}

\begin{theorem}[Planck Depth Formula]
\label{thm:planck-depth}
\[
\nP = 1 + \left\lceil\log_{d+1}\!\left(\frac{\tau_{\mathrm{osc}}}{d\cdot\tP}\right)\right\rceil.
\]
\end{theorem}

\begin{proof}
Solving $d(d+1)^{\nP-1} \geq \tau_{\mathrm{osc}}/\tP$ for $\nP$
yields $\nP - 1 \geq \log_{d+1}(\tau_{\mathrm{osc}}/(d\tP))$; the
minimum integer satisfying this is the given expression. \qed
\end{proof}

\begin{theorem}[LHC Planck Depth]
\label{thm:lhc-planck-depth}
For the LHC bunch clock ($\nu_B = 40.079\,\mathrm{MHz}$,
$\tau_B = 24.951\,\mathrm{ns}$) with $d = 3$ independent detector
subsystems, the Planck depth is $\nP = 60$, reached in physical time
$t_{60} = 60/\nu_B = 1.496\,\mu\mathrm{s}$.
\end{theorem}

\begin{proof}
\begin{align}
\frac{\tau_B}{\tP} &= \frac{24.951 \times 10^{-9}}{5.391 \times 10^{-44}}
= 4.628 \times 10^{35}.
\end{align}
Applying Theorem~\ref{thm:planck-depth} with $d = 3$:
\begin{align}
\nP &= 1 + \left\lceil\log_4\!\left(\frac{4.628 \times 10^{35}}{3}\right)\right\rceil
= 1 + \left\lceil\log_4(1.543 \times 10^{35})\right\rceil \notag\\
&= 1 + \left\lceil\frac{\ln(1.543 \times 10^{35})}{\ln 4}\right\rceil
= 1 + \left\lceil\frac{81.10}{1.386}\right\rceil
= 1 + \lceil 58.51 \rceil = 60.
\end{align}
Verification: $T(60,3) = 3 \cdot 4^{59} = 3 \cdot 3.076 \times 10^{35}
= 9.23 \times 10^{35} > 4.628 \times 10^{35}$ \checkmark. \qed
\end{proof}

\begin{corollary}[Universality of Planck Depth]
For any physically realizable oscillator (frequency in $[\nu_{\min},
\nu_{\max}]$ with $\tP^{-1} > \nu_{\max}$ and $\nu_{\min} > 0$),
the Planck depth at $d = 3$ satisfies $\nP \in [36, 73]$. The
variation spans only 37 ticks despite frequencies varying over 14
orders of magnitude, because $\nP$ depends logarithmically on $\nu$.
\end{corollary}

\begin{table}[h]
\centering
\caption{Planck depth for representative LHC-relevant oscillators ($d=3$).}
\label{tab:planck-depths}
\begin{tabular}{lccr}
\toprule
Oscillator & $\nu$ & $\tau/\tP$ & $\nP$ \\
\midrule
LHC bunch clock & 40.079\,MHz & $4.63 \times 10^{35}$ & \textbf{60} \\
Cs-133 hyperfine & 9.193\,GHz & $2.02 \times 10^{33}$ & 56 \\
H maser & 1.420\,GHz & $1.31 \times 10^{34}$ & 57 \\
Sr optical clock & 429\,THz & $4.32 \times 10^{28}$ & 48 \\
HL-LHC ($d{=}4$) & 40.079\,MHz & $4.63 \times 10^{35}$ & \textbf{52} \\
\bottomrule
\end{tabular}
\end{table}

The HL-LHC row demonstrates the dimensional advantage: adding a fourth
independent detector subsystem (precision timing, $d: 3\to 4$) reduces
$\nP$ from 60 to 52, recovering 8 bunch crossings (200\,ns) of
categorical resolution headroom.

\begin{theorem}[HL-LHC Dimension Advantage]
\label{thm:hllhc-advantage}
For the LHC bunch clock with $d = 4$ (adding the precision timing
detector planned for HL-LHC):
\[
\nP(d=4) = 1 + \left\lceil\log_5\!\left(\frac{4.628 \times 10^{35}}{4}\right)\right\rceil
= 52.
\]
The dimensional upgrade saves $\nP(3) - \nP(4) = 8$ bunch crossings
at fixed oscillator frequency.
\end{theorem}

\begin{proof}
$\log_5(4.628 \times 10^{35}/4) = \log_5(1.157 \times 10^{35})
= \ln(1.157 \times 10^{35})/\ln 5 = 80.5/1.609 = 50.03$;
$\nP = 1 + 51 = 52$. \qed
\end{proof}

% =========================================================
\section{Temporal Programming and Structural Incorruptibility}
\label{sec:temporal}
% =========================================================

\subsection{The Timing Deviation as Sole Datum}

\begin{definition}[$\dP$-Datum]
For a reference oscillator with period $T_{\mathrm{ref}}$ and a
detector signal arriving at time $t_{\mathrm{rec}}(k)$, the
\emph{timing deviation} of the $k$-th signal is
\[
\dP(k) = T_{\mathrm{ref}}(k) - t_{\mathrm{rec}}(k).
\]
In the LHC context, $T_{\mathrm{ref}}(k)$ is the $k$-th bunch crossing
time from the 40\,MHz bunch clock, and $t_{\mathrm{rec}}(k)$ is the
recorded hit time in any detector channel.
\end{definition}

\begin{definition}[Timing Cell]
A \emph{timing cell} $C_i$ is a Borel-measurable subset of the
$\dP$-space $\R$, with associated action $a_i$. The collection
$\{C_i\}$ forms the \emph{cell registry} $\Gamma$, which must
partition the relevant range of $\dP$ values.
\end{definition}

\begin{definition}[Temporal Trigger Program]
A \emph{temporal trigger program} is a tuple $\langle \sigma, \Gamma,
\tau, \phi \rangle$ where $\sigma \in \{\mathrm{IDLE}, \mathrm{COMPILE},
\mathrm{EXECUTE}\}$ is the phase machine state, $\Gamma$ is the cell
registry, $\tau$ is the timing reference channel specification, and
$\phi$ is the action dispatch function mapping cell membership to
trigger decisions.
\end{definition}

\subsection{Structural Incorruptibility}

\begin{theorem}[Structural Incorruptibility of Timing Triggers]
\label{thm:incorruptibility}
A trigger program that conditions entirely on $\dP(k)$ values---with
no content parser applied to signal amplitude, charge, or identity---
has no injection surface: any signal whose timing $\dP(k)$ does not
originate from the LHC bunch crossing will fail to satisfy any
physically calibrated cell condition.
\end{theorem}

\begin{proof}
The cell registry partitions $\dP$-space. A signal from a real
collision produced at the bunch crossing time $T_{\mathrm{ref}}(k)$
satisfies $|\dP(k)| < \delta$ for the bunch-crossing timing jitter
$\delta$ (typically $\lesssim 1\,\mathrm{ns}$ for calibrated
detectors~\citep{atlas_trigger_run2}). A signal from noise, cosmic
rays, or beam halo arrives with a $\dP(k)$ value drawn from a
distribution with mean offset from $T_{\mathrm{ref}}(k)$. Since the
cell conditions $C_i$ are defined by $\dP$ membership alone, and the
cell boundaries are calibrated to the collision timing distribution,
there is no content layer (no message, no encoding, no data packet)
on which an attacker can operate. The only way to satisfy a cell
condition is to produce a signal with the correct timing relative to
$T_{\mathrm{ref}}$, which requires physical presence at the
interaction point. \qed
\end{proof}

\begin{remark}
This theorem identifies the timing trigger as the particle physics
analogue of a structurally incorruptible communication channel: the
``message'' (the physics process) is observed through timing
structure, not transmitted as content. Compare with the Phase-Locked
Messaging framework~\citep{landauer1961irreversibility}, where
security emerges from topological structure rather than cryptographic
secrets.
\end{remark}

\subsection{The Trigger as a Tempus Program}

We write the basic di-muon trigger as a temporal program in outline:

\begin{enumerate}[leftmargin=1.5em,itemsep=0.2em]
\item \textbf{Reference}: The LHC bunch clock at 40.079\,MHz provides
$T_{\mathrm{ref}}(k)$.
\item \textbf{Cells}: $C_{\mathrm{prompt}} = \{|\dP| < 12.5\,\mathrm{ns}\}$;
$C_{\mathrm{background}} = \{12.5 < |\dP| < 100\,\mathrm{ns}\}$;
$C_{\mathrm{noise}} = \{|\dP| > 100\,\mathrm{ns}\}$.
\item \textbf{Composition}: Require two prompt signals in the muon
subsystem within one bunch crossing interval (coincidence).
\item \textbf{Dispatch}: On coincident pair in $C_{\mathrm{prompt}}$,
fire $a = \mathrm{ACCEPT}$; otherwise $a = \mathrm{DISCARD}$.
\end{enumerate}

The composition rule ``two prompt signals within one crossing'' is the
physical trigger condition formalised as a labeled composition of
depth 2 in the muon cell registry.

% =========================================================
\section{The Trigger as Backward Trajectory Completion}
\label{sec:backward}
% =========================================================

\subsection{The Trigger Decision Problem as Trajectory Completion}

The fundamental task of the trigger is: given a pattern of detector
hits (the \emph{endpoint}), determine which physics process produced
it (the \emph{root}). This is precisely backward trajectory completion.

\begin{definition}[Ternary Refinement Hierarchy]
A \emph{ternary refinement hierarchy} of depth $k$ is a partition
structure $(\Omega, \{P_j\}_{j=0}^k, \pi)$ where $|P_0| = 1$,
$|P_{j+1}| = 3|P_j|$, and each block of $P_j$ refines into exactly
three children in $P_{j+1}$. The total leaf count is $N = 3^k$.
\end{definition}

In the trigger context: depth $k$ corresponds to $k$ levels of
detector subsystem hierarchy (L1 calorimeter primitives, L1 muon
system, HLT tracking); each level refines the event classification
into three broad categories (hadronic background, semi-leptonic
candidates, fully leptonic/rare signals).

\begin{definition}[S-Entropy Embedding]
\label{def:s-embed}
The \emph{S-entropy embedding} of a ternary refinement hierarchy of
depth $k$ is the map $\iota_k\colon P_k \to [0,1]^3$ sending each
leaf $p \in P_k$ to the triple $(\Sk(p), \St(p), \Se(p))$ given by
the base-3 expansion of the leaf's three coordinate strings:
\[
\Sk(p) = \sum_{j=1}^{k}\frac{a_j^k(p)}{3^j},\quad
\St(p) = \sum_{j=1}^{k}\frac{a_j^t(p)}{3^j},\quad
\Se(p) = \sum_{j=1}^{k}\frac{a_j^e(p)}{3^j},
\]
where $a_j^c(p) \in \{0,1,2\}$ are the address digits of $p$ at level $j$.
\end{definition}

\begin{theorem}[Trigger State Inflation]
\label{thm:trigger-inflation}
A trigger system with $d$ independent detector subsystems, embedded
via $\iota_k$ and observing $n = k$ levels of hierarchy, can
discriminate $T(n,d) = d(d+1)^{n-1}$ distinct event categories.
\end{theorem}

\begin{proof}
Apply Theorem~\ref{thm:inflation}: the $n = k$ levels of the
hierarchy correspond to $n$ oscillation ticks; the $d$ detector
subsystems correspond to $d$ entropy dimensions. Each labeled
composition of $n$ ticks in $d$ dimensions is a distinct trajectory
through the hierarchy, and distinct trajectories map to distinct
leaves under the embedding $\iota_k$ by the base-3 expansion
(different address strings). \qed
\end{proof}

\subsection{Backward Navigation Complexity}

\begin{theorem}[Backward Navigation Complexity]
\label{thm:backward-complexity}
Backward navigation from any leaf $p_k \in P_k$ to the root $p_0$
via the S-entropy embedding visits exactly $\log_3 N = k$ intermediate
states, where $N = 3^k$, achievable in $\mathcal{O}(\log_3 N)$ steps.
\end{theorem}

\begin{proof}
Each backward step from $p_j$ to $p_{j-1}$ drops the last digit of
the base-3 expansion of $\iota_j(p_j)$, an $\mathcal{O}(1)$ operation.
Starting at $p_k$ and ending at $p_0$ requires $k = \log_3 N$ steps.
\qed
\end{proof}

\begin{theorem}[Collapse Without Virtual States]
\label{thm:collapse}
Backward trajectory completion restricted to intermediate states in
$[0,1]^3$ (physical sub-states only) has worst-case complexity
$\Theta(N)$.
\end{theorem}

\begin{proof}
Without sub-coordinates outside $[0,1]^3$, the backward step from
$p_j$ requires enumerating the children of each candidate parent at
level $j-1$. At level $j$ there are $3^j$ candidate parent blocks,
each with three children. Summing over all levels:
$\sum_{j=0}^{k} 3^j = (3^{k+1}-1)/2 = \Theta(N)$. \qed
\end{proof}

\begin{corollary}
Virtual intermediate states are not optional for efficient trigger
computation: restricting to physically realised intermediate states
increases the computational cost from $\mathcal{O}(\log N)$ to
$\Theta(N)$, making real-time event selection at 40\,MHz impossible.
\end{corollary}

\subsection{Virtual Sub-States in the Trigger}

\begin{definition}[Virtual Sub-State]
A \emph{virtual sub-state} is an intermediate decomposition
$(s_1, s_2, s_3) \in \R^3$ with mean $\bar{s} = s \in [0,1]$
but at least one component $s_i \notin [0,1]$.
\end{definition}

In the trigger context, virtual sub-states arise naturally: when the
HLT reconstructs a track, it uses the Kalman filter equations, which
produce intermediate state estimates that may transiently exceed
physical bounds before converging. These are not artifacts---they are
the mechanism by which the filter achieves its $\mathcal{O}(N_{\rm
det})$ rather than $\mathcal{O}(N_{\rm det}^2)$ complexity.

% =========================================================
\section{Virtual Particles as Virtual Sub-States}
\label{sec:virtual}
% =========================================================

\subsection{Review of Virtual Particles in QFT}

In quantum field theory, a \emph{virtual particle} is an internal
line in a Feynman diagram~\citep{feynman1949space,peskin1995introduction}.
A real (external) particle with mass $m$ satisfies the on-shell
condition $q^2 = q^\mu q_\mu = m^2 c^4$ in the metric convention
$(+,-,-,-)$. A virtual particle has 4-momentum $q^\mu$ satisfying
$q^2 \neq m^2 c^4$; it is \emph{off-shell}. Its contribution to an
amplitude is given by the Feynman propagator:
\begin{equation}
D_F(q) = \frac{i}{q^2 - m^2 c^4 + i\epsilon}.
\end{equation}
The propagator diverges as $q^2 \to m^2$, which is the on-shell limit.
Virtual particles cannot be directly observed---only the external
(on-shell) particles appear as asymptotic states in the
S-matrix~\citep{heisenberg1943smatrix,dyson1949smatrix}.

\subsection{The Formal Identification Theorem}

We now establish the formal correspondence between virtual particles
in QFT and virtual sub-states in the S-entropy framework.

\begin{theorem}[Virtual Particle Identification]
\label{thm:virtual-id}
Virtual particles in quantum field theory are formal virtual sub-states
in the S-entropy framework. Specifically, the following five
correspondences hold:
\begin{enumerate}[leftmargin=1.5em,label=(\roman*),itemsep=0.2em]
\item \textbf{Off-shell $\leftrightarrow$ Outside $[0,1]^3$}: The
on-shell condition $q^2 = m^2 c^4$ corresponds to the physical
constraint $\mathbf{S} \in [0,1]^3$. Off-shell states ($q^2 \neq
m^2 c^4$) correspond to virtual sub-states ($\mathbf{S} \notin
[0,1]^3$).
\item \textbf{Propagator $\leftrightarrow$ Mean-recovery}: The Feynman
propagator $D_F(q) = i/(q^2 - m^2 + i\epsilon)$ encodes the
constraint that virtual momentum sums to give the physical external
momentum. This is the mean-recovery constraint: $(s_1+s_2+s_3)/3 =
s \in [0,1]$.
\item \textbf{Gauge invariance $\leftrightarrow$ Local-Global Decoupling}:
Gauge invariance of the S-matrix means the physical amplitude is
independent of the gauge choice (the local S-values of virtual
contributions). This is precisely the Local-Global Decoupling Theorem:
the global S-value is independent of the local S-values of
subtasks.
\item \textbf{S-matrix observability $\leftrightarrow$ Path Opacity}:
In QFT, only external-state amplitudes are observable; the virtual
intermediate states are unobservable. This is Path Opacity: two
backward trajectories sharing an endpoint are observationally
indistinguishable.
\item \textbf{Unitarity $\leftrightarrow$ Cascade Power}: The
unitarity of the S-matrix ($SS^\dagger = 1$) corresponds to
$\sum_i \catpow_i$ diverging, i.e.\ the cascade of virtual
contributions saturates to complete problem resolution.
\end{enumerate}
\end{theorem}

\begin{proof}
We verify each correspondence.

\textit{(i)}: The physical state space of relativistic particles is
the mass hyperboloid $\mathcal{M}_m = \{q^\mu : q^2 = m^2\}$ in
Minkowski space. The physical state space in the S-entropy framework
is the unit cube $[0,1]^3$. The mapping is: on-shell $\leftrightarrow$
inside the physical domain; off-shell $\leftrightarrow$ outside. This
is a domain identification, not a numerical equality.

\textit{(ii)}: The Feynman propagator at a vertex enforces 4-momentum
conservation: if an external pair $(p_1, p_2)$ produces a virtual
particle of momentum $q = p_1 + p_2$, the propagator weight
$i/(q^2 - m^2 + i\epsilon)$ encodes the constraint that the virtual
momentum sums to give the physical input momentum. Analogously, the
mean-recovery constraint $(s_1+s_2+s_3)/3 = s$ encodes that virtual
sub-components sum (mean) to give the physical global coordinate. Both
are linear constraints on intermediate values.

\textit{(iii)}: In gauge theory, the S-matrix satisfies
$\mathcal{M} = \mathcal{M}(\text{physical states})$ independent of
the gauge parameter $\xi$ in the gauge-fixing term $-(\partial^\mu
A_\mu)^2/(2\xi)$. The Local-Global Decoupling Theorem states: for any
global S-value $\mathcal{S}^*$ and any assignment of local S-values to
subtasks, there exists an expression realising both simultaneously. In
the gauge theory interpretation: for any physical amplitude
$\mathcal{M}$ (global value), the local amplitudes of individual
Feynman diagrams (local S-values) can take any gauge-dependent value
without changing $\mathcal{M}$~\citep{peskin1995introduction}. This is a
statement about freedom of virtual contribution assignment subject
only to the constraint that the sum gives $\mathcal{M}$.

\textit{(iv)}: S-matrix theory~\citep{heisenberg1943smatrix,dyson1949smatrix}
defines physical observables exclusively from the asymptotic states
$|p_1 \cdots p_n\rangle_{\mathrm{in}}$ and $|k_1 \cdots k_m\rangle_
{\mathrm{out}}$. No measurement of any intermediate virtual state is
possible. Path Opacity states: two backward trajectories sharing an
endpoint are indistinguishable by any metric invariant computed at the
endpoint. The endpoint of the backward trajectory is the external state
of the S-matrix; the intermediate virtual sub-states are the internal
virtual particles.

\textit{(v)}: The optical theorem~\citep{cutkosky1960singularities}
gives $2\,\mathrm{Im}\,\mathcal{M}(p \to p) = \sum_X \int d\Phi_X
|\mathcal{M}(p \to X)|^2$, which is the constraint that the sum of
all transition probabilities equals 1 (unitarity). The Cascade Power
formula $\catpow(\Gamma) = 1 - \prod_i(1-\catpow_i)$ saturates to 1
when $\sum_i \catpow_i = \infty$~\citep{borel1909probabilites}. In
QFT, summing over all virtual contributions (all orders of perturbation
theory) gives a unitary result when the perturbation series is
well-defined. \qed
\end{proof}

\subsection{Feynman Diagrams as Labeled Compositions}

\begin{proposition}[Feynman Diagram Count]
\label{prop:feynman-count}
The number of topologically distinct Feynman diagrams at perturbative
order $n$ in a quantum field theory with $d$ field types (particle
species) is bounded below by $T(n,d) = d(d+1)^{n-1}$.
\end{proposition}

\begin{proof}
A Feynman diagram at order $n$ consists of $n$ vertices, each labeled
by one of $d$ interaction terms (field types). The $n$ vertices can be
arranged into an ordered sequence of $k$ groups for any composition
$(c_1, \ldots, c_k)$ of $n$, and each group can be of any of $d$
types. This is precisely a labeled composition of $n$ in $d$
dimensions, and distinct labeled compositions produce topologically
distinct diagrams (different vertex orderings and type assignments
yield different loop structures). \qed
\end{proof}

\begin{remark}
This provides a combinatorial lower bound on the number of Feynman
diagrams. The known growth of QED diagrams as $(2n-1)!! = 1 \cdot 3
\cdot 5 \cdots (2n-1)$ for $n$-loop contributions exceeds
$T(n,3) = 3 \cdot 4^{n-1}$ only for large $n$; the composition-inflation
formula provides the combinatorial floor, while additional symmetry
factors and loop topologies contribute above this floor.
\end{remark}

\subsection{Ward--Takahashi Identities as Local-Global Decoupling}

The Ward--Takahashi identities~\citep{ward1950identity,takahashi1957wardidentity}
are the statement that in QED, the photon propagator and vertex function
satisfy
\begin{equation}
q^\mu \Gamma_\mu(p, p') = S_F^{-1}(p') - S_F^{-1}(p),
\end{equation}
where $\Gamma_\mu$ is the full vertex and $S_F$ is the full fermion
propagator. This identity ensures gauge invariance of the physical
amplitude despite the gauge-dependence of individual diagrams.

\begin{theorem}[Ward-Takahashi as Local-Global Decoupling]
\label{thm:ward-takahashi}
The Ward--Takahashi identity is the gauge-field instantiation of
the Local-Global Decoupling Theorem: it states that the global
physical amplitude is preserved under arbitrary rerouting of
virtual photon momenta (local S-value changes), subject only to
4-momentum conservation (mean-recovery constraint).
\end{theorem}

\begin{proof}
The Ward--Takahashi identity arises from the fact that the gauge
transformation $A_\mu \to A_\mu + \partial_\mu \Lambda$ changes
individual Feynman diagram amplitudes but not the physical S-matrix.
Each diagram amplitude (local S-value) is gauge-dependent; their sum
(global S-value) is not. The Local-Global Decoupling Theorem states
that for any global S-value and any assignment of local S-values to
subtasks, the expression realising both simultaneously exists. The
gauge parameter $\xi$ plays the role of the free local S-value
assignment: varying $\xi$ changes the individual diagram amplitudes
without changing the physical observable. The 4-momentum conservation
at each vertex is the mean-recovery constraint: the sum of incoming
virtual momenta equals the sum of outgoing, exactly as the mean of
the virtual sub-coordinates equals the physical global coordinate. \qed
\end{proof}

% =========================================================
\section{Angular Reformulation of Fundamental Constants}
\label{sec:constants}
% =========================================================

\subsection{The Partition Propagation Bound}

From the partition structure of bounded phase space, the maximum
information propagation rate is finite. For a system with partition
depth $n$ and level spacing $\ell_n$, the propagation speed satisfies
\begin{equation}
c = \sup_n \ell_n \omega_n < \infty,
\end{equation}
where $\omega_n$ is the oscillation frequency at level $n$.
For one complete oscillation cycle ($\Delta\phi = 2\pi$\,rad), the
distance traversed is $\Delta x = 2\pi r_0$ at the characteristic
radius $r_0$. Expressing distance in angular units $\theta = x/r_0$:

\begin{theorem}[Speed of Light in Categorical Angular Units]
\label{thm:c-angular}
\begin{equation}
c = 2\pi\;\mathrm{rad/tick}.
\end{equation}
\end{theorem}

\begin{proof}
One oscillation cycle traverses $\Delta x = 2\pi r_0$ in time
$\Delta t = 1\,\mathrm{tick}$. The speed $c = \Delta x/\Delta t
= 2\pi r_0/\mathrm{tick}$. In angular units (radians, measuring
arc length at $r_0$), $\Delta\theta = \Delta x/r_0 = 2\pi$, so
$c = 2\pi\,\mathrm{rad/tick}$. \qed
\end{proof}

\begin{theorem}[Planck's Constant in Categorical Angular Units]
\label{thm:hbar-angular}
\begin{equation}
\hbar = E_{\mathrm{tick}} / (2\pi),
\end{equation}
where $E_{\mathrm{tick}}$ is the energy carried per oscillation tick.
\end{theorem}

\begin{proof}
From $E = \hbar\omega$ with $\omega = 2\pi/\mathrm{tick}$:
$E_{\mathrm{tick}} = \hbar \cdot 2\pi$. Solving: $\hbar =
E_{\mathrm{tick}}/(2\pi)$. \qed
\end{proof}

\begin{theorem}[Boltzmann's Constant in Categorical Angular Units]
\label{thm:kB-angular}
\begin{equation}
\kB = E_{\mathrm{tick}} / \ln(d+1),
\end{equation}
where $d$ is the dimensionality of the entropy space and $\ln(d+1)$
is the information per composition step.
\end{theorem}

\begin{proof}
Each additional oscillation tick increases the trajectory count by
factor $(d+1)$ (Proposition~\ref{prop:growth}), contributing
$\ln(d+1)$ natural units of entropy. The thermal energy per degree of
freedom is $\kB T$, where $T$ is the tick rate (ticks per second).
Equating the energy per unit of composition entropy: $\kB =
E_{\mathrm{tick}}/\ln(d+1)$. For $d = 3$: $\kB = E_{\mathrm{tick}}/
\ln 4$. \qed
\end{proof}

\begin{theorem}[Planck Time in Categorical Angular Units]
\label{thm:tP-angular}
\[
\tP = \sqrt{\hbar G / c^5} = \frac{\sqrt{E_{\mathrm{tick}} \cdot G}}{(2\pi)^3}.
\]
The Planck time has two components: $(2\pi)^3$ from the angular
geometry of oscillation and $\sqrt{E_{\mathrm{tick}} \cdot G}$ as
the unique irreducible product coupling tick energy to gravitational
coupling.
\end{theorem}

\begin{proof}
Substitute Theorems~\ref{thm:c-angular}--\ref{thm:hbar-angular}:
$\tP = \sqrt{E_{\mathrm{tick}} G / (2\pi \cdot (2\pi)^5)} =
\sqrt{E_{\mathrm{tick}} G}/(2\pi)^3$. \qed
\end{proof}

\begin{corollary}[Irreducibility of $G$]
The gravitational constant $G$ is the unique fundamental constant that
does not reduce to a geometric factor ($2\pi$), a counting quantity
($n$, $d$), or an energy unit ($E_{\mathrm{tick}}$) in the categorical
angular framework. It is the sole external coupling input to the
partition framework.
\end{corollary}

\begin{table}[h]
\centering
\caption{Reduction of fundamental constants to categorical angular units.}
\label{tab:constants}
\begin{tabular}{lll}
\toprule
Constant & SI expression & Categorical form \\
\midrule
$c$ & $2.998 \times 10^8\,\mathrm{m/s}$ & $2\pi\,\mathrm{rad/tick}$ \\
$\hbar$ & $1.055 \times 10^{-34}\,\mathrm{J\cdot s}$ & $E_{\mathrm{tick}}/(2\pi)$ \\
$\kB$ & $1.381 \times 10^{-23}\,\mathrm{J/K}$ & $E_{\mathrm{tick}}/\ln(d{+}1)$ \\
$m_0$ & $\hbar\omega_0/c^2$ & $E_{\mathrm{tick}}\nu_0/(4\pi^2)$ \\
$E_0 = m_0 c^2$ & $m_0 c^2$ & $E_{\mathrm{tick}}\nu_0$ \\
$\tP$ & $\sqrt{\hbar G/c^5}$ & $\sqrt{E_{\mathrm{tick}}G}/(2\pi)^3$ \\
$G$ & $6.674\times 10^{-11}\,\mathrm{m^3 kg^{-1} s^{-2}}$ & \textbf{irreducible} \\
\bottomrule
\end{tabular}
\end{table}

% =========================================================
\section{Discussion}
\label{sec:discussion}
% =========================================================

\subsection{The Current LHC Trigger in Composition-Inflation Terms}

The ATLAS Level-1 trigger has a fixed latency of 2.5\,$\mu$s $= 100$
bunch crossings~\citep{atlas_trigger_run2}. In composition-inflation
terms, this corresponds to depth $n = 100$ in $d = 3$ S-entropy
dimensions, giving a trigger pattern space of $T(100,3) = 3 \cdot
4^{99} \approx 3.0 \times 10^{59}$. The Planck depth is $\nP = 60$;
the trigger operates at $n/\nP = 1.67$, well above the Planck depth.

The current ATLAS trigger menu contains approximately $10^3$ distinct
trigger paths~\citep{atlas_trigger_run2}, corresponding to depth
$n \approx \log_4(10^3/3) + 1 = 6$ --- only 10\% of the Planck
depth. The composition-inflation framework therefore predicts that
$\approx 90\%$ of the available categorical resolution in the L1
timing is unused by the current trigger specification. This suggests
substantial room for improvement in trigger algorithms that exploit
composition structure rather than flat coincidence conditions.

\subsection{Cascade Catalysis and Renormalization}

\begin{theorem}[Renormalization as Cascade Saturation]
\label{thm:rg-saturation}
The Wilsonian renormalization group
flow~\citep{wilson1971renormalization,callan1970broken,symanzik1970small}
corresponds to cascade saturation: the running coupling constant
$g(\mu)$ approaches a fixed point iff the sum of catalytic powers
of the virtual loop corrections diverges.
\end{theorem}

\begin{proof}
The Callan--Symanzik equation $\mu\,\partial g/\partial\mu = \beta(g)$
describes how the coupling $g$ changes with the renormalization scale
$\mu$. A fixed point $g^*$ satisfies $\beta(g^*) = 0$. In the cascade
picture, each loop correction at order $n$ contributes a catalytic
power $\catpow_n = |\beta_n| g^{n-1}$ (proportional to the
loop-correction amplitude). The total catalytic power of the cascade
of loop corrections from scale $\mu_0$ to $\mu$ is
$\catpow_{\mathrm{total}} = 1 - \prod_n(1-\catpow_n)$. The cascade
saturates to power 1 (i.e.\ $g(\mu)$ reaches $g^*$) iff
$\sum_n \catpow_n = \infty$, which by the Borel--Cantelli
analogy~\citep{borel1909probabilites} holds iff the $\beta$-function
series does not decay faster than $1/n$. For asymptotically free
theories ($g^* = 0$, $\beta(g) = -\beta_0 g^2 + \ldots$), the sum
$\sum_n \beta_n g^{n-1}$ diverges as $g \to 0$ from above, confirming
saturation to the trivial fixed point. \qed
\end{proof}

\subsection{Dark Matter and Virtual-State Suppression}

\begin{theorem}[Dark Sector as Virtual-State-Suppressed]
\label{thm:dark-matter}
A particle sector in which backward trajectory completion is restricted
to physical sub-states (Theorem~\ref{thm:collapse}) has $\Theta(N)$
interaction complexity and consequently cannot reach thermal
equilibrium with the visible sector on any finite timescale proportional
to $\log N$.
\end{theorem}

\begin{proof}
If the dark sector is constrained to physical intermediate states
($\mathbf{S} \in [0,1]^3$ at every step), Theorem~\ref{thm:collapse}
gives interaction complexity $\Theta(N)$. The visible sector, using
virtual sub-states, has interaction complexity $\mathcal{O}(\log N)$.
The ratio of dark-to-visible interaction rates is $\Theta(N)/
\mathcal{O}(\log N) = \Theta(N/\log N) \to \infty$. For the
cosmological particle count $N \sim 10^{80}$, the dark sector
interaction rate is suppressed by $10^{80}/\log(10^{80}) \approx
10^{78}$ relative to the visible sector. Gravitational coupling, which
depends only on $G$ (the unique irreducible constant), operates at
the same rate for both sectors, providing the only dark--visible
coupling. \qed
\end{proof}

\begin{remark}
Theorem~\ref{thm:dark-matter} does not constitute a dark matter model
but identifies the categorical structure that would produce the
observed phenomenology: invisible to electromagnetic interactions
(no off-shell photon exchange without virtual sub-states),
gravitationally coupled (through $G$, the irreducible constant),
and thermally decoupled from the visible sector.
\end{remark}

% =========================================================
\section{Conclusion}
\label{sec:conclusion}
% =========================================================

We have established a unified mathematical framework connecting the
real-time event selection problem of LHC trigger systems to categorical
mechanics, composition-inflation, and S-entropy trajectory completion.
The framework produces six principal results:

\textbf{(1) Composition-Inflation (Theorem~\ref{thm:inflation})}: A
trigger with $d$ subsystems and $n$ bunch crossings can discriminate
$T(n,d) = d(d+1)^{n-1}$ event patterns---exponential, not linear,
in $n$.

\textbf{(2) LHC Planck Depth (Theorem~\ref{thm:lhc-planck-depth})}:
$\nP = 60$ for the LHC at $d = 3$; the L1 trigger latency of 100
bunch crossings operates at 1.67 times the Planck depth, placing it
in the super-Planck categorical resolution regime.

\textbf{(3) Structural Incorruptibility (Theorem~\ref{thm:incorruptibility})}:
Timing-only triggers have no injection surface; security is topological,
not cryptographic.

\textbf{(4) Virtual Particle Identification (Theorem~\ref{thm:virtual-id})}:
Virtual particles in QFT are formal virtual sub-states; their five
properties (off-shell, zero net work, path opacity, catalytic power,
necessity) are categorical theorems.

\textbf{(5) Constant Reduction (Table~\ref{tab:constants})}: Six of
seven fundamental constants reduce to tick counts and $2\pi$ geometry;
$G$ is the unique irreducible coupling.

\textbf{(6) HL-LHC Advantage (Theorem~\ref{thm:hllhc-advantage})}:
The precision timing upgrade ($d: 3\to 4$) reduces $\nP$ from 60 to
52, recovering 200\,ns of categorical resolution headroom.

The framework provides formal trigger design principles grounded in
first-principles physics: the optimal trigger depth is $\nP$,
determined by the bunch clock frequency and the number of independent
subsystems; adding subsystems reduces $\nP$ logarithmically; and
the trigger menu is bounded not by engineering constraints alone but
by the composition-inflation state space $T(n,d)$.

\section*{Acknowledgements}
This work was supported by the AIMe Registry for Artificial Intelligence.

\bibliographystyle{plainnat}
\bibliography{references}

\end{document}
